\clearpage
\section{Introduction}
\label{sec:introduction}

Post-training has become a standard stage for enhancing LLMs after pretraining \citep{lai2025survey, cheng2025empowering}, improving instruction-following, reasoning power, alignment with human preferences, and the power of agentic tasks with long-horizon executions.
In particular, recent efforts on LLM post-training mainly focus on Continued Pre-Training (CPT), Supervised Fine-Tuning~(SFT)~\citep{mecklenburg2024injecting} and reinforcement learning~\citep{schulman2017proximal, shao2024deepseekmath,li2026rethinking}, together with dedicated agentic environments ~\citep{glm5team2026glm5vibecodingagentic,dong2026agent} and data pipelines~\citep{nvidia2026nemotron3superopen} for general or domain-specific tasks. 


As LLMs evolve from simple chatbots into agentic workforces, the growing demand for long-context agentic reasoning continues to drive model scaling in pursuit of stronger intelligence. Recent frontier LLMs~\citep{deepseekai2025deepseekv3technicalreport, glm5team2026glm5vibecodingagentic} have pushed model sizes toward the trillion-parameter scale while adopting the hybrid attention mechanism~\citep{yang2025gated,dao2024transformers} or intricate residual connection design~\citep{mhc2025} to improve the efficiency of long-context computation and related optimizations~\citep{chen-etal-2026-dynamic-long, li-etal-2026-lycheecluster}. 
In the meantime, Mixture-of-Experts~(MoE)~\citep{fedus2022switch} architecture introduces a different route for scaling by routing each token to a selected subset of expert networks rather than activating all parameters, per-token computation grows sub-linearly with model size, enhancing reasoning power while preserving memory efficiency.

Most prior studies on large-scale MoE post-training remain tightly coupled with CUDA- or TPU-oriented training infrastructure~\citep{deepseekai2025deepseekv3technicalreport,kimiteam2026kimik2openagentic,glm5team2026glm5vibecodingagentic}, which are built around either Single-Instruction, Multiple-Thread~(SIMT) execution or systolic-array-based accelerator architectures. 
In contrast, post-training trillion-parameter-class LLMs on Single-Instruction, Multiple-Data~(SIMD) hardware remains comparatively underexplored. 
Motivated by this research gap, this report first presents a system-level infrastructure optimization for trillion-parameter-class LLM post-training based on Ascend CloudMatrix384 SuperPOD with Ascend 910C NPUs. 
As one of the recent state-of-the-art LLMs for agentic tasks and reasoning, DeepSeek-V4~\citep{deepseekai2026deepseekv4highlyefficientmilliontoken} reaches top-tier intelligence via intricate sparse attention and residual connection mechanisms. Full-parameter post-training DeepSeek-V4-Pro model on SuperPOD cluster introduces a host of challenges, spanning AscendC kernel design, communication orchestration, memory management, and multi-dimensional parallelism, as illustrated in Figure~\ref{fig:ascend-posttrain-dashboard}.  

Beyond the challenges of training infrastructure, recent LLM post-training efforts have increasingly concentrated on verifiable reasoning domains, particularly mathematical reasoning, coding, and agentic software-engineering tasks. However, many long-tailed industrial decision-making domains remain comparatively under-represented. Operations Research~(OR), for example, is central to real-world optimization problems such as scheduling, transportation, supply-chain management, and resource allocation, yet its integration with LLM training and evaluation is still at an early stage. 
Recent studies on LLM-based OR mainly focus on automatic model formulation and solver-oriented program synthesis, where LLMs translate natural-language problem descriptions into mathematical formulations or executable optimization code. Representative works~\citep{xiao2024chain,zhang2025or} further introduce domain-knowledge retrieval, multi-agent collaboration, and execution-feedback-based refinement into the OR modeling pipeline. However, these approaches remain limited by small-scale benchmarks, brittle semantic-to-structure mapping, and insufficient guarantees of formulation correctness in complex industrial OR scenarios.
Furthermore, existing LLM-based OR studies primarily emphasize agentic-system design, while domain-specific training for OR remains under-explored. This gap becomes more pronounced for trillion-parameter-class models~(e.g., DeepSeek-V4-Pro), where OR-oriented post-training requires not only domain knowledge acquisition via SFT or CPT, but also a dedicated data preparation pipeline, and verifiable feedback from optimization environments. 

In this work, we present full-stack end-to-end training for the DeepSeek-V4 model family based on Ascend SuperPOD with 910C NPUs~\citep{zuo2025serving}. 
We refer to the overall framework as \textbf{\modelname}, which couples full-stack Ascend SuperPOD training optimization with solver-grounded CPT--SFT specialization for OR.
To improve system-level training efficiency, this work first presents a full-stack optimization for SuperPOD-based training infrastructure with both kernel-level optimization and top-level training pipeline orchestration. In particular, we propose AuraKernel, a specialized AscendC kernel optimization agent which automatically optimizes the dominant bottleneck kernels of the DeepSeek-V4 model family. 
By combining the hierarchical optimization strategies, we achieve a 2.93$\times$ Model FLOPs Utilization~(MFU) improvement with the DeepSeek-V4-Pro model. 
Exploiting the optimized system-level training stack on Ascend SuperPOD, we further validate an
OR-oriented post-training workflow on DeepSeek-V4-Flash, covering
solver-grounded data construction, SFT data cleaning, CPT-to-SFT transfer, and
benchmark evaluation. In the CPT-to-SFT transfer evaluation under identical SFT
conditions, direct SFT improves B4O-Feasible from 60.47\% to 65.93\% and B4O-ORGEval from 34.26\% to 48.73\%, while CPT-initialized SFT further raises them
to 71.22\% and 59.39\%, respectively. Aggregating over multiple OR evaluation metrics, our model attains an average accuracy of 71.81\%, outperforming GPT-5.4-Mini by 3.98 points and the base DeepSeek-V4-Flash by 11.27 points. These results show that SFT provides
supervised task alignment, while CPT contributes OR-domain modeling priors that
improve both solver-facing feasibility and structural equivalence, with the
largest gain on B4O-ORGEval dataset. Together, the Pro and Flash post-training validate
two complementary benefits: infrastructure
scalability and domain adaptation. To the best of our knowledge, this is among the first publicly documented full-stack engineering studies of full-parameter post-training for a trillion-parameter-class MoE model on Ascend SuperPOD and it demonstrate the effectiveness of SIMD-based NPUs for large-scale LLM training. 

Overall, the contributions of this work are: 

\begin{itemize}
\item \textbf{Ascend-native system optimization for training complex trillion-parameter MoE LLMs:}
Based on Ascend CloudMatrix384 SuperPOD, we present a full-stack optimization framework for trillion-parameter-scale DeepSeek-V4-Pro training, achieving 34.22\% MFU with 2.93$\times$ improvement over the default baseline. 
In addition to the training acceleration, this work characterizes the key principles of \emph{Ascend-friendly system-level optimization} with jointly optimized parallelism, communication orchestration, memory movement, and kernel execution under the architectural constraints of Ascend SuperPOD. 

\item \textbf{Hardware-aware analysis of MoE, sparse attention, and memory hierarchy on Ascend:}
Through detailed profiling of DeepSeek-V4 training on Ascend, we provide an end-to-end analysis of how model architecture interacts with Ascend kernel execution, communication patterns, memory hierarchy, and distributed orchestration. 
The resulting insights bridge model-level design choices with hardware-level computation efficiency, and comprehensive profiling and bottleneck analysis further motivate the proposed optimization strategies and future research.

\item \textbf{Solver-guided AscendC kernel optimization beyond manual and heuristic optimization:}
We present one of the first attempts to apply Operations Research to Ascend kernel optimization. In contrast to prior kernel-agent approaches that mainly rely on harness construction and loop-level engineering, the proposed AuraKernel introduces OR-guided tiling optimization into the agentic kernel optimization workflow. The AscendC OR solver and dedicated harness design enable accurate tiling selection, long-horizon kernel optimization, and more systematic exploration of the Ascend kernel search space.

\item \textbf{SuperPOD-aware orchestration for complex MoE communication and parallelism:}
We present an end-to-end optimization recipe for Ascend SuperPOD with multi-dimensional parallelism and optimized training orchestration for large-scale MoE architectures.

\item \textbf{Reproducible post-training workflow for domain-specific enhancement:}
We construct an OR-oriented CPT-SFT workflow on DeepSeek-V4-Flash, covering solver-grounded CPT data construction, self-distilled SFT data generation, contract-aware cleaning, and matched benchmark evaluation.

\item \textbf{Empirical evidence for CPT–SFT synergy in structured OR reasoning:}
We show that SFT provides supervised task alignment and improves the original checkpoint, while CPT supplies additional OR-domain modeling priors that further improve both solver-facing feasibility and structural equivalence, with the largest additional gain observed on B4O-ORGEval dataset.

\end{itemize}

